\documentclass[10pt,a4paper]{amsart}
\usepackage[margin=19mm]{geometry}
\usepackage{amsmath,amssymb,mathtools}
\usepackage[scr=rsfs,cal=euler]{mathalfa}
\usepackage{xfrac}
\usepackage[unicode,hidelinks,bookmarks=false]{hyperref}
\newcommand{\ie}{i.e.\ }
\newcommand{\cf}{cf.\ }
\newcommand{\resp}{respectively\ }
\newcommand{\Subsection}{\subsection}
\providecommand{\nobreakdash}{\nobreak\hbox{-}}
\setlength{\emergencystretch}{2em}
\multlinegap0pt
\usepackage{bm}
\usepackage{bbm}
\usepackage{dsfont}
\usepackage{xspace}
\usepackage{cancel}
\usepackage{mathtools}
\usepackage{amsthm}
\makeatletter
\@ifundefined{thm}{\newtheorem{thm}{Thm}}{}
\@ifundefined{cor}{\newtheorem{cor}[thm]{Corollary}}{}
\makeatother
\DeclareMathOperator{\Id}{Id}
\DeclareMathOperator{\Scal}{Scal}
\DeclareMathOperator{\Sym}{Sym}
\newcommand{\deriv}[2]{\frac{\partial #1}{\partial #2}}
\title[Higher Rank Bergman Kernels on Compact Riemann Surfaces]{Higher Rank Bergman Kernels on Compact Riemann Surfaces}
\author{Shin Kim}
\date{Source: arXiv:2505.12241v2 (CC BY 4.0)}
\begin{document}
\maketitle

\noindent\emph{Source and licence.} Higher Rank Bergman Kernels on Compact Riemann Surfaces. Version arXiv:2505.12241v2 (CC BY 4.0), distributed under the Creative Commons Attribution 4.0 International licence, as recorded in the arXiv licence metadata for this exact version and shown on the abstract page https://arxiv.org/abs/2505.12241v2. Full rights notice: licenses/arxiv-2505.12241-CC-BY-4.0.txt.\par\smallskip

\noindent\textit{Editorial scope.} Counted unit: the Cor whose source label is \texttt{cor: HE} in the article (source0.tex lines 2062--2075, source label \texttt{cor: HE}), delivered together with the article's own complete original proof of it (source0.tex lines 2076--2194, one \texttt{proof} environment of 119 lines). No mathematics is written, repaired or re-derived: statement and proof are the article's own text line for line. Required context retained uncounted: none. Every definition, display and label of those lines is retained unchanged. Omitted from this package and not redistributed: 1--2061, 2195--2412. Nothing inside a retained excerpt is omitted or altered: every retained body is delivered exactly as the article's lines are written, so no omission receipt of any kind accompanies this package. Citations: the retained excerpts carry no citation command at all, so no bibliography entry of the article is transcribed and no reference is redistributed. Reference adaptation: none was needed and none was made; every reference inside the delivered unit resolves inside it, so no unresolved reference or citation is produced. Printed slips kept literal: no printed word, symbol, hypothesis or number of the counted statement or of its proof was altered. Layout and class: the article's own class is \texttt{amsart} with packages including amssymb, amsfonts, xy, enumerate, mathrsfs, graphicx, geometry; the wrapper is the lane's pinned \texttt{amsart} editorial wrapper at 10pt on A4, which restates the theorem-like environments the retained text needs and adds the article's own macro definitions that the retained text needs (\texttt{\textbackslash Id}, \texttt{\textbackslash Scal}, \texttt{\textbackslash Sym}, \texttt{\textbackslash deriv}), with extra wrapper packages (bm, bbm, dsfont, xspace, cancel, mathtools). The delivered numbering is the unit's own. Assets: no retained span contains a figure environment, a graphics-inclusion command, a PDF-page inclusion, a table or a TikZ picture; no image file is copied into this package, the package declares no asset dependency and no image file is redistributed with it. Licence: version arXiv:2505.12241v2 of this article is distributed under the Creative Commons Attribution 4.0 International licence, as recorded in the arXiv licence metadata for this exact version and shown on the abstract page https://arxiv.org/abs/2505.12241v2. A full rights notice is delivered at licenses/arxiv-2505.12241-CC-BY-4.0.txt. Characters: every retained excerpt is pure ASCII, so the delivered engine, XeLaTeX with its default Latin Modern text font, reports no missing character.
\begin{cor} \label{cor: HE}
Suppose that $h$ is a Hermitian-Einstein metric. Let $N$ and $p$ be fixed nonnegative integers. Then, there exist smooth functions $b_{1}(x), \dotsc, b_N(x)$, that do not depend on $k$, such that
\begin{equation*}
    B_k(x,x) = \left ( b_{0}(x) \Id_{\Sym^k E}\right )k + \dotsb + \left ( b_{N}(x) \Id_{\Sym^k E} \right ) k^{1-N} + O(k^{-N})
\end{equation*}
where the error term is bounded with respect to the $C^p$ norm. In particular,
\begin{equation*}
b_{0}(x) = c
\end{equation*}
where $c\Id_E = \sqrt{-1} \Lambda F_h$ and 
\begin{equation*}
b_{1}(x) = \frac{1}{2}\Scal_\omega(x).
\end{equation*}
\end{cor}

\begin{proof}

Let $U \subseteq X$ be a sufficiently small coordinate neighborhood and let $\rho \in C_c^\infty(U)$. We know that
\begin{equation*}
\left \|\rho(x) \left ( e^{-\mathfrak{s}^k(\phi(x))} \overline{K_k(x,x)} - \frac{1}{2\pi} \left ( b_{k,0}(x,\overline{x})k + \dotsb + b_{k,N+p}(x,\overline{x}) k^{-N-p+1} \right) \right) \right \|_{C^p,op} = O(k^{-N})
\end{equation*}
and that
\begin{equation*}
    \|b_{k,i}(x,\overline{y})\|_{op(\Sym^k E_x,\Sym^k E_x)} = O(1).
\end{equation*}
We will show that $b_{k,i}(x,\overline{y})$ are of the form $b_{k,i}(x,\overline{y}) = b_i(x,\overline{y}) \Id_{\Sym^k E}$, where $b_i(x,\overline{y})$ is analytic in $x$ and $\overline{y}$, for each $i$. Then, Cauchy's integral formula and the fact that 
\begin{equation*}
    \nabla_0 (b_i(x,\overline{x}) \Id_{\Sym^k E}) = d (b_i)(x,\overline{x}) \otimes \Id_{\Sym^k E}
\end{equation*}
will imply that
\begin{equation*}
    \|\rho(x)b_{k,i}(x,\overline{x})\|_{C^p,op} = O(1).
\end{equation*}


Set 
\begin{equation*}
H(x,z) := h(x,z) = e^{-\psi(x,z)}
\end{equation*}
and write
\begin{equation*}
    \omega(y) = \sqrt{-1} \tilde{g}(y,\overline{y}) dy \wedge d\overline{y}.
\end{equation*}
We know that
\begin{equation*}
\begin{split}
    \tau(x,y,z) = &  - \partial_2(\partial_1 (H) {H}^{-1})(x,z) \\
    & - \frac{1}{2} \partial_2(\partial_1^2(H) {H}^{-1})(x,z)(y-x)  \\
    & + \partial_1(H) {H}^{-1} \partial_2(\partial_1(H) {H}^{-1})(x,z)(y-x) \\
    & + \text{h.d.t.}.
\end{split}
\end{equation*}
On the other hand,
\begin{equation*}
\begin{split}
    \tilde{F}(x,z) & := -\partial_2(\partial_1 (H){H}^{-1})(x,z) \\
    & = \tilde{g}(x,z) \sqrt{-1} \Lambda F(x,z) \\
    & =  c k \tilde{g}(x,z).
\end{split}
\end{equation*}
where $F = F_{h}$ is the curvature form of $h$. In particular, $\tilde{F}(x,z)$ is a scalar function and 
\begin{equation*}
    \begin{split}
        \partial_1 \tilde{F}(x,z) & = -\partial_2(\partial_1^2 (H){H}^{-1})(x,z) + \partial_2 (\partial_1 (H){H}^{-1} \partial_1(H) {H}^{-1})  \\
        & = -\partial_2(\partial_1^2 (H){H}^{-1})(x,z) + 2 \partial_1 (H){H}^{-1} \partial_2 (\partial_1(H) {H}^{-1})  \\
\end{split}
\end{equation*}
is also a scalar function. So, we can write
\begin{equation*}
\begin{split}
    \mathfrak{s}^k(\tau(x,y,z)) =   kf_0(x,z) + kf_1(x,z) (y-x) + \text{h.d.t.}
\end{split}
\end{equation*}
where $f_0(x,z)$ and $f_1(x,z)$ are scalar functions. 



We start with the case when $i = 0$. By our observations from above,
\begin{equation*}
    b_{k,0}(x,z)  = \frac{1}{k} \mathfrak{s}^k(\tau(x,x,z)) \tilde{g}(x,z)^{-1} = c  \Id.
\end{equation*}
Proceeding by induction, suppose that $i > 0$ and that \begin{equation*}
    b_{k,i-1} (x,z) = b_{i-1}(x,z)
\end{equation*}
for some analytic function $b_{i-1}(x,z) \Id_{\Sym^k E}$. If $i = 1$, then
\begin{equation*}
\begin{split}
    (x-y)A_{k,0}(x,y,z) & = k \mathfrak{s}^k(\tau(x,y,z))^{-1} \tilde{g}(y,z) b_{k,0}(x,z) - \Id_{\Sym^k E}  \\
    & = \alpha(x,z) (y-x) + \text{h.d.t.}
\end{split}
\end{equation*}
for some scalar function $\alpha(x,z)$ that does not depend on $k$. Similarly, if $i > 1$, then
\begin{equation*}
\begin{split}
    (x-y) A_{k,i-1}(x,y,z) & = k \mathfrak{s}^k(\tau(x,y,z))^{-1} \tilde{g}(y,z) b_{k,i-1}(x,z)  -  k \mathfrak{s}^k(\tau(x,x,z))^{-1} \tilde{g}(x,z) b_{k,i-1}(x,z)\\
    & = \alpha(x,z) (y-x) + \text{h.d.t.}.
\end{split}
\end{equation*}
So, $A_{k,i-1}(x,x,z) = -\alpha(x,z)$ is a scalar function. Thus,
\begin{equation*}
b_{k,i}(x,z) = \tilde{g}(x,z)^{-1} \deriv{A_{k,i-1}}{z}(x,x,z)
\end{equation*}
is a scalar function and we can write
$$
b_{k,i}(x,z) = b_i(x,z) \Id_{\Sym^k E}
$$
where $b_i(x,z)$ is a function that does not depend on $k$.

Finally, we will compute $b_1(x,\overline{x})$. Since $h$ is Hermitian-Einstein, 
\begin{equation*}
\nabla^{0,1} (\Lambda (F_{\Sym^k h})) = \overline{\partial} (\Lambda (F_{\Sym^k h})) = 0.    
\end{equation*}
Additionally,
\begin{equation*}
\begin{split}
\sqrt{-1} \Lambda (\Delta (F_{ h}))(y) =  & - \tilde{g}^{-1}\partial_1\partial_2(\tilde{g}^{-1}) \tilde{F}(y,\overline{y}) - \tilde{g}^{-1} \partial_1(\tilde{g}^{-1}) \partial_2\tilde{F}(y,\overline{y}) \\
    & + \tilde{g}^{-1} \partial_2(\tilde{g}^{-1}) \eta \tilde{F} (y,\overline{y}) + \tilde{g}^{-2} \eta \partial_2 \tilde{F}(y,\overline{y}) \\
    & -\tilde{g}^{-1} \partial_2(\tilde{g}^{-1}) \partial_1\tilde{F}(y,\overline{y}) - \tilde{g}^{-2} \partial_1\partial_2\tilde{F}(y,\overline{y}) \\
    & - \tilde{g}^{-1} \partial_2(\tilde{g}^{-1}) \tilde{F} \eta(y,\overline{y}) - \tilde{g}^{-2} \partial_2\tilde{F}\eta(y,\overline{y}) \\
    = & \hspace{3pt} \tilde{g}^{-1} \partial_1\partial_2(-\tilde{g}^{-1} \tilde{F}) (y,\overline{y}) \\
    & - \tilde{g}^{-1} \eta \partial_2(-\tilde{g}^{-1} \tilde{F}) (y,\overline{y})\\
    & +  \tilde{g}^{-1} \partial_2(-\tilde{g}^{-1} \tilde{F}) \eta (y,\overline{y}). \\
\end{split}
\end{equation*}
Since $-\tilde{g}^{-1} \tilde{F}(y,\overline{y}) = \sqrt{-1}\Lambda (F_{ h}) = c \Id$, we deduce that $\sqrt{-1} \Lambda(\Delta(F_{h})) = 0$. Then, 
\begin{equation*}
    \sqrt{-1} \Lambda(\Delta(F_{\Sym^k h})) =  \mathfrak{s}^k (\sqrt{-1} \Lambda(\Delta(F_{h}))) = 0.
\end{equation*}
Thus,
\begin{equation*}
 b_{k,1}(x,\overline{x}) = \frac{1}{2} \Scal_\omega(x) \Id_{\Sym^k E}.
\end{equation*}

\end{proof}

\end{document}
